% !TEX program = pdflatex
\documentclass[10pt]{article}
\usepackage[utf8]{inputenc}\usepackage[T1]{fontenc}\usepackage[english]{babel}
\usepackage{lmodern,microtype}\usepackage[margin=.72in]{geometry}
\usepackage{amsmath,amssymb,amsthm,mathtools,booktabs,array,xcolor,enumitem}
\usepackage[most]{tcolorbox}\usepackage{fancyhdr}\usepackage[hidelinks]{hyperref}
\definecolor{CourseBlue}{HTML}{1A3C68}\definecolor{CourseGreen}{HTML}{216E39}
\definecolor{CourseGold}{HTML}{B8860B}\definecolor{CourseRed}{HTML}{9A2C2C}
\definecolor{SoftBlue}{HTML}{F2F6FA}\definecolor{SoftGreen}{HTML}{F1F8F3}
\definecolor{SoftGold}{HTML}{FFF8E8}\definecolor{SoftRed}{HTML}{FFF3F3}
\setlength{\parindent}{0pt}\setlength{\parskip}{.32em}
\setlist[itemize]{leftmargin=1.5em,itemsep=.12em,topsep=.2em}
\pagestyle{fancy}\fancyhf{}\lhead{Stochastic Calculus and Option Pricing}
\rhead{Lecture 14 Notes}\cfoot{\thepage}\renewcommand{\headrulewidth}{.3pt}
\newcommand{\E}{\mathbb E}\newcommand{\Q}{\mathbb Q}\newcommand{\cF}{\mathcal F}
\newcommand{\dd}{\,\mathrm d}\newcommand{\PhiN}{\Phi}
\newtcolorbox{boardexample}[1]{enhanced,breakable,colback=SoftBlue,colframe=CourseBlue,boxrule=.55pt,arc=1mm,title={Example --- #1},fonttitle=\bfseries}
\newtcolorbox{studentexercise}[1]{enhanced,breakable,colback=SoftGold,colframe=CourseGold,boxrule=.55pt,arc=1mm,title={Exercise --- #1},fonttitle=\bfseries}
\newtcolorbox{shortsolution}[1][]{enhanced,breakable,colback=SoftGreen,colframe=CourseGreen,boxrule=.55pt,arc=1mm,title={Short solution#1},fonttitle=\bfseries}
\newtcolorbox{keyidea}[1]{enhanced,breakable,colback=white,colframe=CourseGreen,boxrule=.7pt,arc=1mm,title={#1},fonttitle=\bfseries}
\newtcolorbox{warningbox}[1]{enhanced,breakable,colback=SoftRed,colframe=CourseRed,boxrule=.65pt,arc=1mm,title={#1},fonttitle=\bfseries}
\title{\textbf{Pricing Toolbox: A Guided Refresher}\\\large Lecture 14 Notes --- It\^o, Girsanov, Path Dependence, and Early Exercise}
\author{}\date{}
\begin{document}\maketitle\vspace{-1.4em}
\begin{keyidea}{The four questions}
\[
\begin{array}{c|c}
\text{Question}&\text{Tool}\\\hline
\text{How do we simplify a nonlinear Brownian quantity?}&\text{It\^o formula}\\
\text{How do we remove a drift?}&\text{Girsanov}\\
\text{What information must a payoff remember?}&\text{state augmentation}\\
\text{When should the holder stop?}&\text{dynamic programming}
\end{array}
\]
The methods change a representation while preserving the pricing problem.
\end{keyidea}

\section{It\^o calculus turns integrals into identities}
For \(f(t,B_t)\),
\[
\dd f(t,B_t)=f_t\dd t+f_x\dd B_t+\frac12f_{xx}\dd t.
\]
Taking \(f(x)=e^x\) gives
\[
\dd e^{B_t}=e^{B_t}\dd B_t+\frac12e^{B_t}\dd t,
\]
hence
\[
\boxed{\int_0^T e^{B_s}\dd B_s=e^{B_T}-1-\frac12\int_0^T e^{B_s}\dd s.}
\]
\begin{boardexample}{the stochastic exponential}
\[
\mathcal E_t(\theta)=\exp\left(\theta B_t-\frac12\theta^2t\right)
\]
satisfies \(\dd\mathcal E_t=\theta\mathcal E_t\dd B_t\). The compensating
quadratic-variation term is why it is a martingale rather than merely an
exponential function of Brownian motion.
\end{boardexample}
\begin{studentexercise}{It\^o correction}
Why is \(\exp(\theta B_t)\) not a martingale for \(\theta\ne0\)?
\end{studentexercise}
\begin{shortsolution}
It\^o gives a positive drift
\(\frac12\theta^2\exp(\theta B_t)\dd t\). Subtracting
\(\frac12\theta^2t\) in the exponent cancels that drift.
\end{shortsolution}

\section{Girsanov removes drift by changing measure}
If \(W^\Q_t=W_t+\int_0^t\lambda_s\dd s\) is Brownian motion under a new
measure \(\Q\), then a process with drift under the old measure can become a
driftless Brownian motion under \(\Q\). The density is a stochastic
exponential, subject to an integrability condition.
\begin{keyidea}{The purpose of the change}
Changing measure reweights paths; it does not change observed path values.
For pricing, choose the measure under which discounted traded assets are
martingales. For hitting problems, choose the measure under which the
underlying Brownian calculation is simple.
\end{keyidea}
\begin{boardexample}{first hitting time}
For \(X_t=\mu t+B_t\), the event of hitting level \(a>0\) can be related by
Girsanov to a driftless Brownian hitting event. Reflection then yields the
inverse-Gaussian density
\[
f_{\tau_a}(t)=\frac{a}{\sqrt{2\pi t^3}}
\exp\left(-\frac{(a-\mu t)^2}{2t}\right),\qquad t>0.
\]
\end{boardexample}

\section{Path-dependent claims need an enlarged state}
For an arithmetic-average Asian call,
\[
\left(\frac1T\int_0^TS_u\dd u-K\right)^+,
\]
define
\[
I_t=\int_0^tS_u\dd u,\qquad \dd I_t=S_t\dd t.
\]
The stock alone is not Markov for this payoff; the state is
\((t,S_t,I_t)\). Under Black--Scholes, the value solves
\[
v_t+rs v_s+s v_i+\frac12\sigma^2s^2v_{ss}-rv=0,
\]
with terminal payoff \((I/T-K)^+\).

\begin{warningbox}{Simulation has two errors here}
A Monte Carlo Asian estimator has sampling error and numerical integration
error for the time integral. A geometric Asian option is often a useful
control variate because its price is tractable while it is strongly
correlated with the arithmetic payoff.
\end{warningbox}
\begin{studentexercise}{a barrier-monitoring bias}
Why can a discretely monitored up-and-in option be underpriced relative to
continuous monitoring?
\end{studentexercise}
\begin{shortsolution}
A path may cross the barrier between observation dates and return below it
before the next observation. The discrete estimator misses that activation;
refining the grid reduces the bias.
\end{shortsolution}

\section{Barrier events combine maxima and measure changes}
For an up-and-in one-touch with barrier \(H>S_0\), the payoff is
\[
\mathbf 1_{\{\sup_{0\leq u\leq T}S_u\geq H\}}.
\]
With \(X_t=\log(S_t/S_0)=\mu t+\sigma B_t\), the barrier becomes
\(a=\log(H/S_0)\). Girsanov handles drift and reflection handles the maximum.
The price is the discounted hitting probability.

\section{American options are stopping decisions}
In a binomial model with gross return \(R\) and risk-neutral probability
\(p\),
\[
\boxed{V_t=\max\left\{g(S_t),
\frac{pV_{t+1}^u+(1-p)V_{t+1}^d}{R}\right\}.}
\]
The first term is exercise value and the second is continuation value.
\begin{boardexample}{a three-step American put}
Let \(S_0=100\), \(u=1.1\), \(d=.9\), \(R=1.02\), \(K=100\). Then
\(p=.6\). At node \(S_2=81\),
\[
\text{exercise}=19,\qquad
\text{continuation}=\frac{.6(10.9)+.4(27.1)}{1.02}=17.04.
\]
Exercise is optimal. Working backward gives approximately
\[
V_0^{\rm Eur}=4.59,\qquad V_0^{\rm Am}=4.91.
\]
The difference is the early-exercise premium.
\end{boardexample}

\begin{studentexercise}{the correct operation}
At a node, a put has exercise value \(8\) and continuation value \(9.2\).
What is its American value and action?
\end{studentexercise}
\begin{shortsolution}
The value is \(\max\{8,9.2\}=9.2\). Continue; exercise and continuation
are alternative policies, not cash flows to add.
\end{shortsolution}

\section{Numerical habits complete the toolbox}
\begin{keyidea}{First diagnostic by contract feature}
\begin{center}
\begin{tabular}{p{.30\textwidth}|p{.52\textwidth}}
\toprule
Problem&First check\\\midrule
Asian Monte Carlo&integration grid and sampling error\\
Barrier simulation&monitoring bias and bridge correction\\
American tree&time-step refinement and exercise region\\
PDE or local-volatility grid&boundary placement and mesh convergence\\
\bottomrule
\end{tabular}
\end{center}
\end{keyidea}
\begin{warningbox}{A number is not yet a result}
Before trusting a numerical price, apply a relevant refinement or
no-arbitrage check. A formula can be correct while its discretization,
simulation, or implementation is not.
\end{warningbox}
\begin{keyidea}{Final synthesis}
The payoff tells us what information must be carried forward: quadratic
variation, a probability density, a state variable, a path maximum, or a
stopping rule.
\end{keyidea}
\end{document}
